\subsection{代数不相交扩张}

定义5. --- 设 L 是域 K 的一个扩域，E 和 F 是 L 的两个子扩张。若对于 E（相应地，F）的每个在 K 上代数自由的子集 A（相应地，B），A 与 B 不相交且 A ∪ B 在 K 上代数自由，则称 E 与 F（在 K 上）代数不相交，并称 E 在 K 上与 F 代数不相交。

注记. --- 1) 若 E 是 L 的一个子扩张，且在 K 上是代数的，则它与 L 的每个子扩张 F 都是代数不相交的。一个 K 的扩张是代数的，当且仅当它与自身代数不相交。

\begin{enumerate}
\def\labelenumi{\arabic{enumi})}
\setcounter{enumi}{1}
\item
  可能发生 E 在 K 上与 F 代数不相交、但在 K 的某个子域 \(K_{0}\) 上并非如此的情形。* 例如，C 在 R 上与自身代数不相交，但在 Q 上并非如此。*
\item
  显然，如果 E 在 K 上与 F 代数不相交，那么当 E 和 F 被视为 L 的子扩张时，这一性质在它们被视为 \(\mathbf{K}(\mathbf{E} \cup \mathbf{F})\) 的子扩张时同样成立，反之亦然。
\end{enumerate}

命题11. --- 若E与F在K上代数不相交，则E ∩ F在K上是代数的。

这由定义5得出。

命题12. --- 设L是域K的一个扩域，E、F是L的子扩域。则以下条件等价：

\begin{enumerate}
\def\labelenumi{\alph{enumi})}
\item
  E 与 F 代数不相交；
\item
  存在 E 在 K 上的一个超越基，它在 F 上是代数自由的。
\item
  E 的每个在 K 上代数自由的子集都在 F 上代数自由。我们引入以下条件：
\end{enumerate}

\(b^{\prime})\) 存在 \(\mathbf{F}\) 在 \(\mathbf{K}\) 上的一个超越基，它在 \(\mathbf{E}\) 上代数自由；

\(c^{\prime})\) \(\mathbf{F}\) 的每个在 \(\mathbf{K}\) 上代数自由的子集都在 \(\mathbf{E}\) 上代数自由。

\(a) \Rightarrow b')\) ：设 E 与 F 代数不相交。设 B（相应地，C）是 E（相应地，F）在 K 上的一个超越基。则 \(B \cap C = \emptyset\) 且 \(B \cup C\) 在 K 上代数自由，因此 C 在 \(K(B)\) 上代数自由（V，第107页，命题4）；由于 E 在 \(K(B)\) 上是代数的，V 第108页的命题6表明 C 在 E 上代数自由。

\(b') \Rightarrow c)\) ：假设存在 F 在 K 上的一个超越基 C，且 C 在 E 上是代数自由的。设 A 是 E 的一个在 K 上代数自由的子集。则 C 在 K(A) 上是代数自由的，因此 A 在 K(C) 上是代数自由的（V，第107页，命题4），从而在 F 上也是代数自由的（V，第108页，命题6），因为 F 在 K(C) 上是代数的。

\(c) \Rightarrow a)\) 这由命题4（V，第107页）立即得出。

蕴含关系 \(a) \Rightarrow b) \Rightarrow c') \Rightarrow a)\) 可以用同样的方式证明。

推论. --- 设 E 与 F 在 K 上代数不相交。设 E' 为 E 在 L 中的相对代数闭包，F' 为 F 在 L 中的相对代数闭包（V，第19页）。则 E' 与 F' 在 K 上代数不相交。

设 B 是 E 在 K 上的一个超越基；它也是 \(E'\) 在 K 上的一个超越基。由于 E 在 K 上与 F 代数不相交，B 在 F 上代数自由，从而在 \(F'\) 上代数自由（V，第108页，命题6）；于是可以应用命题12。

命题 13. --- 设 L 是域 K 的一个扩张，且 E、F 是 L 的两个子扩张。

\begin{enumerate}
\def\labelenumi{\alph{enumi})}
\tightlist
\item
  我们有 \(\operatorname{tr} \cdot \deg_{\mathbb{F}} \mathrm{F}(\mathrm{E}) \leqslant \operatorname{tr} \cdot \deg_{\mathbb{K}} \mathrm{E}\) 。当 \(\mathrm{E}\) 与 \(\mathrm{F}\) 在 \(\mathbf{K}\) 上代数不相交时， \(\mathrm{E}\) 在 \(\mathbf{K}\) 上的每个超越基都是
\end{enumerate}

\(F(E)\) 在 F 上的一个超越基，并且我们有 tr. \(\deg_{F}F(E) = \text{tr. } \deg_{K}E\) 。反之，这个等式蕴含 E 与 F 在 K 上代数不相交，当 tr. \(\deg_{K}E\) 有限时。

\begin{enumerate}
\def\labelenumi{\alph{enumi})}
\setcounter{enumi}{1}
\item
  我们有 \(\operatorname{tr} \cdot \deg_{\mathbf{K}} \mathrm{K}(\mathrm{E} \cup \mathrm{F}) \leqslant \operatorname{tr} \cdot \deg_{\mathbf{K}} \mathrm{E} + \operatorname{tr} \cdot \deg_{\mathbf{K}} \mathrm{F}\) 。当 \(\mathrm{E}\) 与 \(\mathrm{F}\) 在 \(\mathbf{K}\) 上代数不相交时，我们有 \(\operatorname{tr} \cdot \deg_{\mathbf{K}} \mathrm{K}(\mathrm{E} \cup \mathrm{F}) = \operatorname{tr} \cdot \deg_{\mathbf{K}} \mathrm{E} + \operatorname{tr} \cdot \deg_{\mathbf{K}} \mathrm{F}\) 。反之，这个等式蕴含 \(\mathrm{E}\) 与 \(\mathrm{F}\) 在 \(\mathbf{K}\) 上代数不相交，当 \(\mathrm{E}\) 与 \(\mathrm{F}\) 在 \(\mathbf{K}\) 上超越次数有限时。
\item
  设 B 是 E 在 K 上的一个超越基；则 E 在 K(B) 上是代数的，且 V 第18页的推论2表明 F(E) 在 F(K(B)) = F(B) 上是代数的。由定理2（V，第109页），B 包含 F(E) 在 F 上的一个超越基；当 E 在 K 上与 F 代数不相交时，B 在 F 上是代数自由的（命题12），此时它就是 F(E) 在 F 上的一个超越基。a) 的前三个断言由此得出。现设 E 在 K 上的超越次数有限，且等于 F(E) 在 F 上的超越次数；由于 F(E) 在 F(B) 上是代数的且 Card B = tr . deg\_F F(E)，V 第110页的推论1表明 B 在 F 上是代数自由的，因此 E 在 K 上与 F 代数不相交（命题12）。
\item
  我们有 \(\mathbf{K}(\mathbf{E} \cup \mathbf{F}) = \mathbf{F}(\mathbf{E})\) ，从而 V 第 110 页的推论蕴含等式：
\end{enumerate}

\[
\operatorname{tr} \cdot \deg_ {K} K (E \cup F) = \operatorname{tr} \cdot \deg_ {F} F (E) + \operatorname{tr} \cdot \deg_ {K} F.
\]

现在 \(b)\) 立即由 \(a)\) 及此等式得出。

命题14. --- 设 L 是域 K 的一个扩域，E 和 F 是 L 的两个子扩张，且 B 是 E 在 K 上的一个超越基。为使 E 与 F 在 K 上代数不相交，充要条件是 K(B) 与 F 在 K 上线性不相交。

为使 E 与 F 在 K 上代数不相交，充要条件是 B 在 F 上代数自由（命题12），即 B 中元素的单项式在 F 上线性无关。由于这些单项式构成向量 K-空间 K{[}B{]} 的一组基，这等价于说 K{[}B{]} 与 F 在 K 上线性不相交。最后，由于 K(B) 是 K{[}B{]} 的分式域，V 第14页的命题6表明 K{[}B{]} 与 F 线性不相交当且仅当 K(B) 与 F 线性不相交。

推论1. --- 若 E 与 F 线性不相交，则 E 在 K 上与 F 代数不相交。反之，若 E 是 K 的一个纯扩张，且在 K 上与 F 代数不相交，则 E 与 F 在 K 上线性不相交。

推论2. --- K 的每个纯扩张与 K 的每个代数扩张都线性不相交；特别地，K 在 K 的每个纯扩张中都是相对代数闭的。

